\documentclass[12pt]{article}

%%%%%%%%%%%%%%%%%%%%%%%%%%%%
%%%%%%%%%%%%%%%%%%%%%%%%%%%%
% Load in packages
\usepackage{amsmath}
\usepackage{float}
\usepackage{amssymb}
\usepackage{hyperref}
\usepackage{graphicx}
\usepackage{enumitem}
\usepackage[top=1in, bottom=1in, left=1in, right=1in]{geometry}

%%%%%%%%%%%%%%%%%%%%%%%%%%%%
%%%%%%%%%%%%%%%%%%%%%%%%%%%%

\begin{document}

\begin{center}
\Large Chapter 14 Practice Problems

\medskip

\normalsize Elements of Microeconomics (discussion section 1)

\medskip

\small Rudy Kretschmar
\end{center}

\section*{Question 1}
\begin{enumerate}[label = (\alph*)]
    \item Draw a graph of cost curves (ATC, AVC, MC) and marginal revenue in a perfectly competitive market in the following situations and explain why they look the way they do:
\begin{enumerate}[label = (\roman*)]
    \item The firm should shutdown in the short run and should exit in the long run
    \item The firm should remain open in the short run but should exit in the long run
\end{enumerate}
    \item In each of cases (i) and (ii) above, draw and shade the profits or losses for the firm in that situation.
\end{enumerate}

\section*{Question 2}
Consider a firm in a perfectly competitive market which produces donuts. The firm has costs of production shown in the following table:

\begin{table}[H]
    \centering
    \begin{tabular}{ccccccc}
    \hline
        Q & Total FC & Total VC & AFC & AVC & ATC & MC \\
        \hline
        \hline
         0 & 50 & 0 & & & &  \\
         1 & 50 & 25 & & & & \\
         2 & 50 & 35 & & & &  \\
         3 & 50 & 55 & & & &  \\
         4 & 50 & 90 & & & &  \\
         5 & 50 & 150 & & & & \\
         \hline
    \end{tabular}
    \caption{Costs of production for a donut firm}
    \label{tab:placeholder}
\end{table}

\begin{enumerate}[label = (\alph*)]
    \item Fill in the missing values in the table
    \item The market price for a donut is \$35. Graph the ATC, AVC, MC, MR/P.
    \item  Should the firm remain open in the short run? Should the firm exit the market in the long run? Use the shutdown and exit conditions to defend your assertions.
\end{enumerate}


\section*{Question 3}
Explain why the short run shutdown condition and long run enter/exit conditions are different. What is the key factor considered in the long run that is not considered in the short run?

\end{document}
